\begin{afterword}
\summaryhead

Posted on April 1st, the post is a mock news story. A team led by the Nobel laureate Santiago Ram\'on y Cajal, with Leeuwenhoek and possibly Imhotep, announces that the brain is a network of tiny cells, revealed by Golgi's stain. Its statement gathers the evidence of centuries: Alcmaeon and Broca on the brain as the seat of reason, Nemesius's claim that the faculties sit in the ventricles, which leaves no reason for an intricate brain, Galvani and Helmholtz on the electrical nature of nerves, and Babbage's engine as a model of thought built from small devices.

The team concludes that finding an enormously complicated material system where the mind should be shows that Spinoza was right and Descartes wrong: mind and body are of one substance. With Darwin's account of how such an organ could arise, the evidence weighs against theories, all religions included, that make the mind fundamental. A full account of thought is promised but not yet delivered. The researchers are no longer available for comment.

\responsehead

Judged as satire, the post does what the author proposed for Amazing Breakthrough Day a week earlier. It reports an old and understandable discovery as news, and the form makes its point: that the brain is an intricate material system where the mind should be is a large fact, easy to overlook because it was established slowly. The history is correct nearly everywhere I could check it: Cajal's prize, Golgi's stain, Alcmaeon, Nemesius's ventricles and the ``too earthy'' brain tissue, Helmholtz's measurement of the nerve signal, the conjecture linking Imhotep to the Edwin Smith Papyrus. The argument against Nemesius is fair: a mind housed in empty cavities gives no reason to expect an intricate tissue around them.

The release's philosophy is less accurate than its history. One substance is Spinoza's thesis, but Spinoza held that thought and extension have nothing in common and do not act on each other; thought is as basic as matter. The release goes on, with Darwin, to the conclusion that intelligence is ``ontologically non-fundamental'', which is not Spinoza's view. The finding fits Spinoza only in part.

Its conclusions are hedged (``seems to indicate'', ``strongly weighs against'', ``the promise, though not yet the realization''). One clause claims more than the rest: the evidence weighs against ``all religions ever invented.'' That assumes every religion makes the mind fundamental, and the release discusses none. And the evidence did not decide the question for everyone who knew it best. John Eccles, who shared a Nobel Prize for work on the synapse, remained a dualist.

In short: an April 1st story that does what the author proposed, with accurate history, and a claim to vindicate Spinoza that fits only part of Spinoza's view.
\end{afterword}
